% =============================================================
%  CC3104 - Aprendizaje por Refuerzo | Laboratorio 6
%  Gradiente de politica: REINFORCE con linea base vs Actor-Critic
%  Plantilla uvglab (version standalone, preambulo inlineado)
%  Compilar: pdflatex lab6_cc3104.tex  (dos veces, por \pageref{LastPage})
%
%  CAMBIOS RESPECTO AL PREAMBULO DEL LAB 4:
%    1. - colores 'congestion', 'estante', 'rutasegura' (eran del mapa TikZ)
%    2. - \usepackage{tikz}  (no hay mapa en este lab; si lo necesitas,
%         descomenta la linea marcada)
%    3. + \usepackage{subcaption}  (las 4 graficas del Task 2 van en grid)
%    4. + \graphicspath{{figs/}{./}}  (para las figuras de la entrega final)
%    5. + \usepackage{amsfonts} y \DeclareMathOperator{\Var}  (varianza,
%         sesgo y esperanzas aparecen por todo el Task 1.3)
%    6. + entorno 'prediccion' (caja con marco punteado) para marcar las
%         predicciones del Task 1 que hay que contrastar en el Task 3.1
%    7. + \usepackage{csquotes} para el resumen del paper (Task 3.4)
%  El resto del preambulo es identico al del Lab 4.
% =============================================================
\documentclass[11pt,letterpaper]{article}

\usepackage[utf8]{inputenc}
\usepackage[T1]{fontenc}
% babel-spanish: si lo tenes instalado (Overleaf lo tiene), se carga solo.
\IfFileExists{spanish.ldf}{\usepackage[spanish,es-tabla]{babel}}{}
\usepackage[margin=2.5cm,top=3cm,bottom=2.8cm]{geometry}
\usepackage{amsmath,amssymb,amsfonts}
\usepackage{xcolor}
\usepackage{fancyhdr}
\usepackage{tcolorbox}
\usepackage{booktabs}
\usepackage{graphicx}
\usepackage{subcaption}
\usepackage{lastpage}
\usepackage{enumitem}
\usepackage{csquotes}
\usepackage[hidelinks]{hyperref}
\usepackage{listings}
% \usepackage{tikz}   % <- descomentar si agregas diagramas

\tcbuselibrary{skins,breakable}

\graphicspath{{figs/}{./}}

\DeclareMathOperator{\Var}{Var}
\DeclareMathOperator{\Sesgo}{Sesgo}
\newcommand{\Esp}{\mathbb{E}}

% ---------- Paleta UVG ----------
\definecolor{uvgverde}{RGB}{0,104,56}
\definecolor{uvgverdeclaro}{RGB}{230,243,235}
\definecolor{uvggris}{RGB}{90,90,90}

% ---------- Encabezado y pie ----------
\pagestyle{fancy}
\fancyhf{}
\fancyhead[L]{\small\textcolor{uvgverde}{\textbf{CC3104 -- Aprendizaje por Refuerzo}}}
\fancyhead[R]{\small Christian Echeverría (221441)\\[-2pt]\small Mario Betancourt (23440)}
\fancyfoot[L]{\small\textcolor{uvggris}{UVG -- Facultad de Ingeniería}}
\fancyfoot[R]{\small\textcolor{uvggris}{Pág. \thepage\ de \pageref{LastPage}}}
\renewcommand{\headrulewidth}{0.8pt}
\renewcommand{\footrulewidth}{0.4pt}
\setlength{\headheight}{26pt}

% ---------- Titulos de seccion ----------
\usepackage{titlesec}
\titleformat{\section}{\Large\bfseries\color{uvgverde}}{\thesection.}{0.6em}{}
\titleformat{\subsection}{\large\bfseries\color{uvgverde}}{\thesubsection}{0.6em}{}
\titlespacing*{\section}{0pt}{1.4em}{0.8em}
\titlespacing*{\subsection}{0pt}{1.1em}{0.6em}

% ---------- Cajas ----------
\newtcolorbox{enunciado}{
  breakable, enhanced,
  colback=white, colframe=uvgverde, boxrule=0.9pt,
  arc=2pt, left=8pt, right=8pt, top=6pt, bottom=6pt,
  title=\textbf{Enunciado}, fonttitle=\small,
  coltitle=white, colbacktitle=uvgverde
}

\newtcolorbox{respuesta}{
  breakable, enhanced,
  colback=uvgverdeclaro, colframe=uvgverde, boxrule=0.9pt,
  arc=2pt, left=8pt, right=8pt, top=6pt, bottom=6pt,
  title=\textbf{Respuesta}, fonttitle=\small,
  coltitle=white, colbacktitle=uvgverde
}

\newtcolorbox{observacion}{
  breakable, enhanced,
  colback=white, colframe=uvggris, boxrule=0.6pt,
  arc=2pt, left=8pt, right=8pt, top=6pt, bottom=6pt,
  title=\textbf{Observación}, fonttitle=\small,
  coltitle=white, colbacktitle=uvggris
}

% Caja para las predicciones del Task 1 que se contrastan en el Task 3.1
\newtcolorbox{prediccion}{
  breakable, enhanced,
  colback=white, colframe=uvgverde, boxrule=0.9pt,
  frame style={dash pattern=on 4pt off 3pt},
  arc=2pt, left=8pt, right=8pt, top=6pt, bottom=6pt,
  title=\textbf{Predicción (se contrasta en Task 3.1)}, fonttitle=\small,
  coltitle=white, colbacktitle=uvgverde
}

\newtcolorbox{pendiente}{
  breakable, enhanced,
  colback=white, colframe=uvggris, boxrule=0.6pt,
  frame style={dash pattern=on 4pt off 3pt},
  arc=2pt, left=8pt, right=8pt, top=6pt, bottom=6pt,
  title=\textbf{Pendiente}, fonttitle=\small,
  coltitle=white, colbacktitle=uvggris
}

\newcommand{\lead}[1]{\textbf{#1}}

% ---------- Estilo de codigo ----------
\definecolor{codegris}{RGB}{245,245,245}
\lstdefinestyle{uvgpython}{
  language=Python,
  backgroundcolor=\color{codegris},
  basicstyle=\ttfamily\footnotesize,
  keywordstyle=\color{uvgverde}\bfseries,
  commentstyle=\color{uvggris}\itshape,
  stringstyle=\color{teal},
  numbers=left, numbersep=8pt, numberstyle=\tiny\color{uvggris},
  frame=leftline, framerule=2pt, rulecolor=\color{uvgverde},
  showstringspaces=false, breaklines=true, tabsize=4, xleftmargin=12pt
}
\lstset{style=uvgpython}

% =============================================================
\begin{document}

% ---------- Bloque de titulo compacto ----------
\begin{center}
  {\Large\bfseries\color{uvgverde} Laboratorio 6}\\[4pt]
  {\large REINFORCE con línea base vs.\ Actor-Critic en control continuo}\\[6pt]
  {\small Universidad del Valle de Guatemala --- Facultad de Ingeniería}\\[2pt]
  {\small CC3104 -- Aprendizaje por Refuerzo}\\[6pt]
  {\small Christian Echeverría (221441) \quad\textbullet\quad Mario Betancourt (23440)}\\[2pt]
  {\small \today}
\end{center}
\vspace{0.6em}
\hrule
\vspace{1.2em}

\begin{observacion}
\lead{Contexto del caso.} Una empresa de robótica de asistencia médica desarrolla un
exoesqueleto para rehabilitación de movilidad en pacientes con lesión de rodilla. El
controlador debe asistir el movimiento de la pierna de forma suave y eficiente,
minimizando el esfuerzo del motor. Antes de pasar a hardware, se valida si los métodos
de gradiente de política son apropiados para este dominio de control continuo, usando
\texttt{LunarLanderContinuous-v2} como proxy del problema real.
\end{observacion}

% =============================================================
\section{Task 1}
% =============================================================

\begin{enunciado}
Respondan las siguientes preguntas con argumentación técnica rigurosa antes de
implementar nada.
\end{enunciado}

\subsection{Inadecuación de Q-Learning tabular y DQN}

\begin{enunciado}
El entorno de simulación es \texttt{LunarLanderContinuous-v2} de Gymnasium, con espacio
de acción continuo de dos dimensiones que representa la fuerza de dos propulsores.
Argumenten formalmente por qué Q-Learning tabular y DQN son inapropiados para este
entorno. El argumento debe mencionar explícitamente el \textbf{espacio de acción}, el
\textbf{operador \textit{argmax}} y la \textbf{representación de la política}.
\end{enunciado}

\begin{respuesta}
\lead{Espacio de acción.} En \texttt{LunarLanderContinuous-v2} el espacio de acción es
$\mathcal{A} = [-1,1]^2 \subset \mathbb{R}^2$, donde $a^{(1)}$ regula el propulsor
principal y $a^{(2)}$ los laterales. Es un conjunto \emph{no numerable}: tiene la
cardinalidad del continuo. El espacio de estados $\mathcal{S} \subset \mathbb{R}^8$
también lo es.

\medskip
\lead{Por qué falla Q-Learning tabular.} El método mantiene una tabla indexada por pares
$(s,a)$ y aplica
\[
  Q(S_t,A_t) \leftarrow Q(S_t,A_t) + \alpha\Big[R_{t+1} + \gamma \max_{a}Q(S_{t+1},a) - Q(S_t,A_t)\Big].
\]
Con $\mathcal{S}\times\mathcal{A}$ no numerable la tabla no es representable en memoria
finita. Pero el problema de fondo no es el almacenamiento: es que la tabla \emph{no
generaliza}. Cada entrada se actualiza de forma independiente, y la probabilidad de
volver a visitar exactamente el mismo par $(s,a)$ es cero (es un conjunto de medida nula
bajo cualquier política con densidad). Cada entrada recibiría a lo sumo una
actualización, de modo que el estimador nunca acumula estadística y el promedio muestral
que justifica la convergencia de Q-Learning jamás se forma. Ni con memoria infinita
aprendería.

\medskip
\lead{Por qué falla DQN.} DQN resuelve el lado del \emph{estado}: sustituye la tabla por
$Q_w(s,\cdot)$, una red que generaliza entre estados parecidos. No resuelve el lado de la
\emph{acción}, por dos razones distintas.

\emph{(i) Arquitectura.} La red de DQN emite un vector de $|\mathcal{A}|$ salidas, una
por acción, precisamente para poder tomar el máximo con un solo paso hacia adelante. Esa
arquitectura presupone $\mathcal{A}$ finito y discreto; con $\mathcal{A}$ continuo el
número de cabezas de salida no está definido.

\emph{(ii) El operador \textit{argmax}.} Es el obstáculo esencial, y aparece en dos
lugares:
\[
  \pi(s) \;=\; \arg\max_{a \in \mathcal{A}} Q(s,a),
  \qquad
  y_t \;=\; R_{t+1} + \gamma \max_{a' \in \mathcal{A}} Q_{w^-}(S_{t+1}, a').
\]
Con $\mathcal{A}$ discreto ambos se resuelven enumerando. Con $\mathcal{A}$ continuo cada
uno es un problema de optimización global sobre $[-1,1]^2$, y $Q_w$ es una red neuronal:
no es cóncava en $a$, no tiene solución en forma cerrada, y admite múltiples óptimos
locales. Además el máximo no se necesita una vez, sino en \emph{cada} transición y para
\emph{cada} muestra del minilote objetivo. Con minilotes de 64 y $10^6$ pasos de
entrenamiento, eso son del orden de $6\times 10^7$ optimizaciones no convexas anidadas
dentro del bucle de aprendizaje. El costo es prohibitivo y, peor aún, cualquier error del
maximizador interno se propaga como sesgo sistemático al objetivo de \textit{bootstrapping}.

\medskip
\lead{Representación de la política.} Esta es la diferencia estructural. En Q-Learning y
DQN la política es \emph{implícita y derivada}: no existe como objeto propio, sino que se
extrae de $Q$ mediante la maximización anterior. De ahí se siguen dos limitaciones. Es
\emph{determinista por construcción} —la única estocasticidad disponible es
$\varepsilon$-greedy, que inyecta ruido uniforme y arbitrario, no una exploración
modulada por el estado— y su evaluación cuesta una optimización completa. Los métodos de
gradiente de política representan $\pi_\theta(a\mid s)$ de forma \emph{explícita y
paramétrica}: la red emite directamente los parámetros de una distribución sobre
$\mathcal{A}$, muestrear una acción cuesta $O(1)$, la exploración es intrínseca y su
magnitud ($\sigma$) es aprendible, y el óptimo se busca por ascenso de gradiente en
$\theta$ en lugar de por maximización sobre $a$. La maximización sobre el espacio de
acción desaparece del algoritmo.

\medskip
\lead{Por qué la discretización no lo salva.} La objeción obvia es discretizar
$\mathcal{A}$ en $k$ niveles por dimensión. Con $d=2$ y $k=11$ dan 121 acciones, que
parece manejable. Falla por cuatro motivos:

\begin{enumerate}[leftmargin=1.4em,itemsep=2pt]
  \item \emph{Crecimiento exponencial.} El número de acciones es $k^d$. El problema real
        —un exoesqueleto con varias articulaciones actuadas— tiene $d$ del orden de 6 o
        7; con $k=11$ eso son $1.9\times 10^7$ salidas. El proxy es engañosamente barato.
  \item \emph{Pérdida de resolución.} El error de cuantización está acotado por
        $\tfrac{2}{k-1}$ por dimensión, y el aterrizaje suave exige control fino del
        empuje cerca del contacto. Aumentar $k$ para recuperar precisión es exactamente
        lo que dispara el punto anterior.
  \item \emph{Destrucción de la estructura métrica.} $\mathcal{A}$ es un espacio métrico:
        empujes cercanos producen efectos cercanos. Al discretizar, cada nivel se vuelve
        una etiqueta simbólica sin relación con sus vecinos, de modo que lo aprendido
        sobre un nivel de empuje no se transfiere al contiguo. Se descarta información
        del problema.
  \item \emph{Conflicto con el requisito de suavidad.} Anticipando la sección 1.4: una
        política sobre acciones cuantizadas solo puede cambiar el empuje en saltos de
        tamaño $\tfrac{2}{k-1}$. Para un dispositivo sujeto a la pierna de un paciente,
        una discretización gruesa \emph{introduce} la brusquedad que hay que evitar.
\end{enumerate}
\end{respuesta}

\subsection{Gradiente del logaritmo de la política Gaussiana}

\begin{enunciado}
Para \texttt{LunarLanderContinuous-v2} la política se parametriza como
$\pi_\theta(a \mid s) = \mathcal{N}(\mu_\theta(s), \sigma^2 I)$, donde $\mu_\theta(s)$ es
la salida de una red neuronal. Expliquen cómo se calcula
$\nabla_\theta \ln \pi_\theta(A_t \mid S_t)$ para esta parametrización específica.
Desarrollen la expresión analítica del gradiente del logaritmo de la densidad Gaussiana
respecto a $\theta$, identificando qué parte depende de $\theta$ y qué parte no.
\end{enunciado}

\begin{respuesta}
\lead{Densidad y logaritmo.} Para $a \in \mathbb{R}^d$ con covarianza isotrópica
$\sigma^2 I$:
\[
  \ln \pi_\theta(a \mid s) \;=\; -\frac{1}{2\sigma^{2}}\,\lVert a - \mu_\theta(s) \rVert^{2}
  \;-\; \frac{d}{2}\ln\!\left(2\pi\sigma^{2}\right)
\]

\medskip
Conviene separar el término logarítmico para exhibir la dependencia en $\sigma$:
\[
  \ln \pi_\theta(a \mid s) \;=\; -\frac{1}{2\sigma^{2}}\,\lVert a - \mu_\theta(s) \rVert^{2}
  \;-\; d\ln\sigma \;-\; \frac{d}{2}\ln(2\pi), \qquad d = 2.
\]

\medskip
\lead{Qué depende de $\theta$ y qué no.} Distinguimos tres bloques:

\begin{enumerate}[leftmargin=1.4em,itemsep=2pt]
  \item $\tfrac{d}{2}\ln(2\pi)$ es una constante de normalización pura: no depende ni de
        $\theta$ ni de $\sigma$ ni de los datos. Su gradiente es nulo y en la
        implementación puede omitirse sin alterar nada.
  \item $\mu_\theta(s)$ es lo \emph{único} que depende de los pesos de la red. Toda la
        dependencia en $\theta$ entra por el término cuadrático.
  \item $\sigma$ merece atención aparte. En la parametrización clásica $\sigma$ es un
        hiperparámetro fijo y entonces $-d\ln\sigma$ también es constante. Pero la
        especificación del Task 2 exige que $\sigma$ sea un \emph{parámetro aprendible}
        (independiente del estado, inicializado en $0.5$). Bajo esa especificación
        $\theta = (\theta_\mu, \rho)$ con $\rho = \ln\sigma$ —se parametriza el
        logaritmo para garantizar $\sigma > 0$ sin restricciones— y el término
        $-d\ln\sigma = -d\rho$ \emph{sí} contribuye al gradiente. Ambos casos se
        desarrollan abajo.
\end{enumerate}

\medskip
\lead{Gradiente respecto a los pesos de la red.} Sea
$J(s) = \partial \mu_\theta(s)/\partial\theta_\mu \in \mathbb{R}^{d\times|\theta_\mu|}$ el
jacobiano que entrega la retropropagación. Derivando el término cuadrático por regla de
la cadena:
\[
  \nabla_{\theta_\mu} \ln \pi_\theta(a \mid s)
  \;=\; \frac{1}{\sigma^{2}}\, J(s)^{\!\top}\big(a - \mu_\theta(s)\big).
\]

\medskip
\lead{Gradiente respecto a la desviación estándar.} Con $\sigma = e^{\rho}$:
\[
  \frac{\partial}{\partial \rho}\ln \pi_\theta(a\mid s)
  \;=\; \frac{\lVert a - \mu_\theta(s)\rVert^{2}}{\sigma^{2}} \;-\; d.
\]

\medskip
\lead{Interpretación.} El factor $(a-\mu_\theta(s))/\sigma^2$ es el residuo entre la
acción efectivamente muestreada y la media propuesta, escalado por la precisión. Es la
dirección en el espacio de acciones que \emph{aumenta} la verosimilitud de $A_t$; en el
estimador completo $\hat{A}_t \nabla_\theta \ln\pi_\theta$ es la ventaja la que decide el
signo: si $\hat{A}_t > 0$ la media se desplaza hacia la acción tomada, si $\hat{A}_t < 0$
se aleja de ella. El gradiente respecto a $\rho$ tiene esperanza nula en el $\sigma$
correcto, ya que $\Esp\lVert a-\mu\rVert^2 = d\sigma^2$: aumenta la dispersión cuando las
acciones muestreadas caen sistemáticamente más lejos de lo que $\sigma$ predice y la
reduce en caso contrario. Ese es el mecanismo por el cual la entropía cae durante el
entrenamiento, y es exactamente lo que se mide en el Task 3.3.

\medskip
\lead{Tres consecuencias que importan para la implementación.}

\emph{(a) La ventaja es constante respecto a $\theta$.} En
$\nabla_\theta J(\theta) = \Esp[\hat{A}_t\,\nabla_\theta \ln\pi_\theta(A_t\mid S_t)]$, tanto
$A_t$ (ya muestreada) como $\hat{A}_t$ se tratan como constantes al derivar. Si el grafo
de cómputo deja que el gradiente fluya hacia el crítico a través de $\hat{A}_t$, el
estimador deja de ser el gradiente de política. En PyTorch esto obliga a un
\texttt{.detach()} sobre la ventaja.

\emph{(b) El factor $1/\sigma^2$ es una fuente de inestabilidad.} La magnitud del
gradiente crece como $\sigma^{-2}$. Si la política colapsa hacia el determinismo, la
norma $\lVert\nabla_\theta L\rVert$ se dispara. Es una predicción directa sobre lo que
debería observarse en el Task 3.2.

\emph{(c) La activación $\tanh$ introduce saturación y desajuste de soporte.} Con
$\mu_\theta(s) = \tanh(f_\theta(s))$, el jacobiano lleva el factor
$1-\tanh^2(f_\theta(s))$, que se anula cuando la preactivación se satura: una media
empujada hacia los extremos del rango deja de recibir gradiente. Además la Gaussiana
tiene soporte en todo $\mathbb{R}^d$ aunque su media esté acotada, de modo que con
$\sigma = 0.5$ una fracción no despreciable de las acciones muestreadas cae fuera de
$[-1,1]$ y el entorno las recorta. La acción ejecutada deja entonces de coincidir con la
acción cuya densidad se usó en el estimador, lo que introduce un sesgo que la formulación
teórica no contempla.
\end{respuesta}

\subsection{Sesgo y varianza: REINFORCE con línea base vs.\ Actor-Critic}

\begin{enunciado}
Comparen formalmente REINFORCE con línea base y Actor-Critic en términos de sesgo y
varianza del estimador del gradiente. Para cada algoritmo identifiquen: qué usa como
estimador de la ventaja $\hat{A}_t$, qué componente introduce sesgo y qué componente
introduce varianza. Predigan cuál algoritmo esperan que converja más rápido en
\texttt{LunarLanderContinuous-v2} y justifiquen esa predicción.
\end{enunciado}

\begin{respuesta}
\lead{Estimador de la ventaja en cada algoritmo.}
\begin{center}
\begin{tabular}{lll}
\toprule
\textbf{Algoritmo} & $\hat{A}_t$ & \textbf{Bootstrapping} \\
\midrule
REINFORCE con línea base & $G_t - \hat{V}_w(S_t)$ & No \\
Actor-Critic (TD(0))     & $R_{t+1} + \gamma \hat{V}_w(S_{t+1}) - \hat{V}_w(S_t)$ & Sí \\
\bottomrule
\end{tabular}
\end{center}

\medskip
\lead{REINFORCE con línea base: por qué es insesgado.} La línea base no introduce sesgo
\emph{cualquiera que sea} su calidad, siempre que no dependa de la acción. Para todo
$b(S_t)$:
\begin{align*}
  \Esp_{a\sim\pi_\theta}\!\big[b(S_t)\,\nabla_\theta \ln\pi_\theta(a\mid S_t)\big]
  &= b(S_t)\int_{\mathcal{A}}\pi_\theta(a\mid S_t)\,
     \frac{\nabla_\theta \pi_\theta(a\mid S_t)}{\pi_\theta(a\mid S_t)}\,da \\[2pt]
  &= b(S_t)\,\nabla_\theta \int_{\mathcal{A}}\pi_\theta(a\mid S_t)\,da
   \;=\; b(S_t)\,\nabla_\theta 1 \;=\; 0.
\end{align*}
El punto es fuerte y conviene subrayarlo: si la red de línea base está mal entrenada, el
estimador pierde eficacia reductora de varianza pero \emph{sigue apuntando en la
dirección correcta en esperanza}. No puede desviar el aprendizaje.

\medskip
\lead{REINFORCE con línea base: de dónde viene la varianza.} $G_t = \sum_{k\ge 0}\gamma^k
R_{t+k+1}$ es una suma Monte Carlo sobre la trayectoria restante completa. Acumula dos
fuentes de aleatoriedad en cada uno de los pasos que faltan: el muestreo de la Gaussiana
y la estocasticidad del entorno. A grandes rasgos la varianza crece con el horizonte
efectivo $\tfrac{1}{1-\gamma}$. En \texttt{LunarLanderContinuous-v2} los episodios llegan
a 1000 pasos y los retornos van de $\approx -400$ a $\approx +300$ según el episodio
termine en choque o en aterrizaje: la dispersión del estimador es enorme. A esto se suma
un problema de eficiencia distinto de la varianza: al requerir el episodio completo, el
algoritmo produce \emph{una sola actualización por episodio}, es decir alrededor de 1000
pasos de optimización en toda la corrida.

\medskip
\lead{Actor-Critic: de dónde viene el sesgo.} El error TD
$\delta_t = R_{t+1}+\gamma\hat{V}_w(S_{t+1})-\hat{V}_w(S_t)$ satisface
$\Esp[\delta_t \mid S_t=s, A_t=a] = A^\pi(s,a)$ \emph{únicamente} si
$\hat{V}_w \equiv V^\pi$. Fuera de ese caso el error de aproximación
$\epsilon(s) = \hat{V}_w(s)-V^\pi(s)$ entra directamente en la dirección del gradiente:
al inicio del entrenamiento $\hat{V}_w$ es esencialmente ruido y el actor se actualiza
hacia un objetivo equivocado. Dos agravantes propios de esta especificación: el crítico
se entrena sobre datos generados por una política que cambia, de modo que persigue un
objetivo no estacionario; y el enunciado obliga a que el crítico \emph{comparta el cuerpo
de la red con el actor}, así que los gradientes del actor modifican las características
sobre las que el crítico predice y viceversa. Ese acoplamiento hace que la escala relativa
de ambas pérdidas se vuelva un hiperparámetro crítico.

\medskip
\lead{Actor-Critic: por qué la varianza es baja.} $\delta_t$ involucra una única
recompensa y dos evaluaciones de la red de valor. Su varianza es $O(1)$ respecto al
horizonte, frente al $O\!\big(\tfrac{1}{1-\gamma}\big)$ de $G_t$. Además admite una
actualización por paso, lo que multiplica por dos o tres órdenes de magnitud el número de
pasos de optimización a igual presupuesto de episodios.

\medskip
\lead{Resumen del compromiso.}
\begin{center}
\begin{tabular}{lcc}
\toprule
 & \textbf{Sesgo} & \textbf{Varianza} \\
\midrule
REINFORCE con línea base & Nulo & Alta \\
Actor-Critic             & Positivo & Baja \\
\bottomrule
\end{tabular}
\end{center}
\end{respuesta}

\begin{prediccion}
\lead{P1 (velocidad de convergencia).} Actor-Critic alcanzará una recompensa media móvil
de 20 episodios igual o superior a $0$ antes del episodio $600$; REINFORCE con línea base
no alcanzará ese umbral dentro de los 1000 episodios de la corrida.

\emph{Justificación.} El factor dominante no es el sesgo sino el conteo de
actualizaciones. Con episodios de varios cientos de pasos, Actor-Critic realiza del orden
de $10^5$ actualizaciones en 1000 episodios frente a las $10^3$ de REINFORCE. A eso se
suma que la recompensa de \texttt{LunarLanderContinuous-v2} es densa y con
\textit{shaping}, condición bajo la cual el \textit{bootstrapping} de TD(0) resulta
informativo desde temprano y su sesgo se amortiza rápido.

\medskip
\lead{P2 (norma del gradiente).} La norma media $\lVert\nabla_\theta L\rVert$ de REINFORCE
será mayor que la de Actor-Critic, y su dispersión entre episodios también, porque
$|G_t - \hat{V}_w|$ es de orden $10^2$ mientras que $|\delta_t|$ es de orden $10^0$.

\medskip
\lead{P3 (entropía).} Partiendo de $\sigma_0 = 0.5$ y $d=2$, la entropía inicial es
$S[\pi_\theta] = \tfrac{d}{2}\ln(2\pi e\sigma^2) \approx 1.45$ nats. Se predice que
decrece de forma monótona en ambos algoritmos y que cae más rápido en Actor-Critic, por
el mismo argumento de conteo de actualizaciones de P1.

\medskip
\lead{Modo de fallo previsto para P1.} Se registra por anticipado la vía por la que P1
podría refutarse: Actor-Critic con TD(0), tronco compartido y sin término de entropía ni
normalización de la ventaja es propenso al colapso prematuro de $\sigma$. Si eso ocurre,
la política se vuelve determinista antes de haber encontrado una región buena, se estanca
en un óptimo local y REINFORCE —lento pero insesgado— podría terminar por encima. En ese
escenario P1 quedaría refutada y P3 confirmada de manera que \emph{explicaría} la
refutación de P1, no que la contradiría.
\end{prediccion}

\subsection{Restricción de suavidad temporal en el exoesqueleto real}

\begin{enunciado}
El entorno de exoesqueleto real tiene una restricción que
\texttt{LunarLanderContinuous-v2} no tiene: las acciones deben ser suaves en el tiempo
para no causar movimientos bruscos que dañen al paciente. Argumenten cómo modificarían
la función de recompensa y la parametrización de la política para incorporar esa
restricción. ¿Cambiaría eso la elección entre REINFORCE y Actor-Critic?
\end{enunciado}

\begin{respuesta}
\lead{Modificación de la recompensa.} Penalización sobre la diferencia de acciones
consecutivas:
\[
  R'_t \;=\; R_t \;-\; \lambda \,\lVert a_t - a_{t-1} \rVert^{2}
\]

con $\lambda > 0$. Si además interesa penalizar la aceleración del comando y no solo su
velocidad de cambio, puede añadirse un término de tercer orden (\textit{jerk})
$-\lambda_2\lVert a_t - 2a_{t-1} + a_{t-2}\rVert^2$, al costo de ampliar más el estado.

\medskip
\lead{Consecuencia sobre la propiedad de Markov.} Este es el punto que la modificación
obliga a tratar y que suele pasarse por alto. Al depender $R'_t$ de $a_{t-1}$, la
recompensa deja de ser función de $(s_t,a_t)$ y el proceso deja de ser un MDP sobre
$\mathcal{S}$. Hay que ampliar el estado:
\[
  \tilde{s}_t \;=\; (s_t,\, a_{t-1}) \;\in\; \mathbb{R}^{8+2}.
\]
Sin esa ampliación $V^\pi(s)$ estaría promediando sobre una variable no observada y tanto
la ecuación de Bellman como las garantías de ambos algoritmos dejarían de aplicar. La
red pasa a recibir 10 entradas.

\medskip
\lead{Calibración de $\lambda$ y riesgo de \textit{reward hacking}.} La penalización tiene
un óptimo degenerado: la acción nula constante la anula por completo. Si $\lambda$ domina
sobre la magnitud típica de $R_t$, la política óptima es no moverse, que es un
exoesqueleto perfectamente suave y perfectamente inútil. Criterio operativo: fijar
$\lambda$ de modo que $\lambda\lVert\Delta a\rVert^2$ represente una fracción pequeña
—del orden del 5 al 10\%— de $|R_t|$ típico, y verificar empíricamente que la política
resultante sigue completando la tarea antes de aceptar el valor.

\medskip
\lead{Modificación de la parametrización de la política.} La penalización en la recompensa
es una restricción \emph{blanda}: el agente puede violarla si le conviene. Para un
dispositivo sujeto a la pierna de un paciente eso es insuficiente, y conviene trasladar
parte de la restricción a la parametrización, donde se vuelve estructural.

\emph{(a) Parametrización incremental con límite de tasa.} En lugar de emitir la acción
absoluta, la red emite un incremento acotado:
\[
  a_t = \mathrm{clip}\big(a_{t-1} + \Delta_t,\,-1,\,1\big),
  \qquad \Delta_t \sim \mathcal{N}\big(\delta_{\max}\tanh(f_\theta(\tilde{s}_t)),\,\sigma^2 I\big).
\]
Esto impone una cota dura sobre $\lVert a_t - a_{t-1}\rVert$ por diseño, no por
incentivo: una garantía de seguridad verificable sin necesidad de confiar en el
aprendizaje. El costo es que introduce un integrador —los errores se acumulan— y que
requiere $a_{t-1}$ en el estado, cosa que de todos modos ya exigía la recompensa.

\emph{(b) Tratar $\sigma$ como el problema, no como un detalle.} Es el punto menos obvio y
probablemente el más importante. Aunque $\mu_\theta$ fuese perfectamente suave, con ruido
independiente en cada paso las acciones consecutivas difieren en
$\Esp\lVert a_t - a_{t-1}\rVert^2 = 2d\sigma^2$, es decir un orden de $1.0$ con
$\sigma=0.5$ sobre un rango de $2$. \emph{La política estocástica es en sí misma la
principal fuente de brusquedad}, con independencia de lo que aprenda la red. Dos remedios:
sustituir el ruido independiente por ruido temporalmente correlacionado (proceso de
Ornstein-Uhlenbeck, como hace DDPG), o desplegar en hardware únicamente la media
determinista $\mu_\theta(\tilde{s})$ y reservar la estocasticidad para el entrenamiento en
simulación.

\emph{(c) Filtrado o memoria.} Un filtro pasa-bajos a la salida, o una política recurrente
que mantenga estado interno, logran suavidad a costa de introducir retardo de fase, lo
que en un lazo de asistencia física puede degradar la sincronía con el paciente.

\medskip
\lead{¿Cambia la elección entre REINFORCE y Actor-Critic?} Sí, y refuerza el argumento de
la sección 1.3 en lugar de matizarlo. La razón es de asignación de crédito. El término
$-\lambda\lVert a_t - a_{t-1}\rVert^2$ es una señal \emph{densa y estrictamente local}:
depende solo de dos acciones consecutivas y su responsable es un único paso. TD(0) la
atribuye exactamente donde ocurre. En cambio el retorno Monte Carlo de REINFORCE reparte
esa penalización sobre todas las acciones anteriores del episodio, de modo que cientos de
decisiones suaves reciben crédito negativo por una sacudida que ninguna de ellas causó.
Eso agrega varianza sin agregar información: la modificación empeora específicamente el
punto débil de REINFORCE.

\medskip
\lead{Salvedad.} Dicho eso, la conclusión honesta es más limitada de lo que parece.
Ninguno de los dos algoritmos es desplegable en hardware: ambos son \textit{on-policy},
ineficientes en muestras, y sin región de confianza una sola actualización grande puede
producir una política peligrosa —precisamente el problema que motivó PPO y que se examina
en el Task 3.2. La pregunta relevante para el sistema real no es cuál de los dos elegir,
sino cuál sirve mejor como \emph{punto de partida} en simulación antes de escalar a un
método con control del tamaño de paso. Ese argumento se retoma en el Task 3.5 con la
evidencia experimental.
\end{respuesta}

% =============================================================
\section{Task 2}
% =============================================================

\subsection{Arquitectura y protocolo experimental}

\begin{respuesta}
% madagascar
Ambos algoritmos se implementaron desde cero con PyTorch. El Actor recibe un estado de
ocho componentes y usa dos capas ocultas de 64 neuronas con ReLU. La cabeza de media
aplica $\tanh$, y un único parámetro escalar $\log\sigma$, inicializado en
$\log(0.5)$, se expande a las dos dimensiones; por tanto, la covarianza es
$\sigma^2I$. Durante entrenamiento se muestrea $a_t\sim\mathcal{N}(\mu_\theta(s_t),
\sigma^2I)$ y se recorta a $[-1,1]^2$ antes de enviarla al entorno. La entropía
reportada es la de esta Gaussiana previa al recorte.

REINFORCE usa una red de valor independiente, también $64$--$64$--$1$, ajustada con
$\operatorname{MSE}(\hat V(S_t),G_t)$. Su pérdida del Actor es
$-\sum_t\log\pi_\theta(A_t\mid S_t)[G_t-\hat V(S_t)]$, con la ventaja separada del
grafo. Actor-Critic comparte las capas ocultas y agrega una cabeza escalar de valor. En
cada transición calcula
$\delta_t=R_{t+1}+\gamma\hat V(S_{t+1})-\hat V(S_t)$ y minimiza
$-\log\pi_\theta(A_t\mid S_t)\delta_t+\tfrac12\delta_t^2$. En una truncación por
tiempo se conserva el \textit{bootstrap}; solo una terminación real lo anula.

Se ejecutaron 1000 episodios por método con semilla 42 y $\gamma=0.99$. Las tasas de
aprendizaje fueron $10^{-3}$ para el Actor y la línea base de REINFORCE, y
$3\times10^{-4}$ para la red compartida de Actor-Critic. Gymnasium 1.2.2 retiró el
identificador \texttt{LunarLanderContinuous-v2}; por ello la ejecución local utilizó su
sucesor oficial \texttt{LunarLanderContinuous-v3}, conservando el espacio de estados de
ocho dimensiones y el espacio de acciones continuo $[-1,1]^2$. El notebook intenta
primero \texttt{v2} para seguir siendo compatible con instalaciones antiguas.

La norma de REINFORCE corresponde a su actualización episódica. En Actor-Critic se
calcula la norma del gradiente exclusivo de la pérdida del Actor en cada transición y se
reporta su promedio episódico; se mide antes de sumar la pérdida del Critic para que el
tronco compartido no contamine la métrica. Esta diferencia de granularidad obliga a
interpretar las escalas de forma cualitativa y no como una comparación de magnitudes
idénticamente normalizadas.
\end{respuesta}

\subsection{Resultados registrados}

\begin{respuesta}
La tabla resume los arreglos completos guardados en
\texttt{results/lab6\_metrics.npz}. El cruce indica el primer episodio en que la media
móvil de 20 recompensas alcanzó cero.

\begin{center}
\resizebox{\textwidth}{!}{%
\begin{tabular}{lrrrrrrrr}
\toprule
\textbf{Algoritmo} & \textbf{Primeros 100} & \textbf{Últimos 100} & \textbf{Mejor}
& \textbf{Cruce} & \textbf{$\|\nabla L_A\|$ med.} & \textbf{$\|\nabla L_A\|$ máx.}
& \textbf{$S_1$} & \textbf{$S_{1000}$} \\
\midrule
REINFORCE      & $-259.80$ & $14.73$   & $96.69$ & 778 & $2287.58$ & $30709.21$ & $1.4516$ & $1.4369$ \\
Actor-Critic   & $-446.09$ & $-147.63$ & $18.31$ & --- & $38.76$   & $515.61$   & $1.4514$ & $1.8816$ \\
\bottomrule
\end{tabular}}
\end{center}

REINFORCE mejoró desde $-259.80$ en los primeros 100 episodios hasta $14.73$ en los
últimos 100; su media móvil terminó en $14.89$ y alcanzó un máximo de $39.01$ en el
episodio 991. Actor-Critic también mejoró respecto de su comienzo, pero se estabilizó en
una solución inferior: su mejor media móvil fue $-125.71$ en el episodio 968 y terminó
en $-141.91$. Ninguno resolvió el entorno bajo el criterio usual de recompensa 200.
\end{respuesta}

\begin{figure}[htbp]
  \centering
  \includegraphics[width=0.96\textwidth]{curvas_aprendizaje.png}
  \caption{Recompensa total con media móvil de 20 episodios. REINFORCE cruza cero en el
  episodio 778; Actor-Critic no alcanza el umbral durante la corrida.}
  \label{fig:curvas}
\end{figure}

\begin{figure}[htbp]
  \centering
  \includegraphics[width=0.96\textwidth]{norma_gradiente.png}
  \caption{Norma del gradiente del Actor por episodio en escala logarítmica.}
  \label{fig:gradientes}
\end{figure}

\begin{figure}[htbp]
  \centering
  \includegraphics[width=0.96\textwidth]{entropia.png}
  \caption{Entropía de la política Gaussiana. La trayectoria no es monótonamente
  decreciente en ninguno de los algoritmos.}
  \label{fig:entropia}
\end{figure}

\begin{figure}[htbp]
  \centering
  \includegraphics[width=0.96\textwidth]{episodio_render.png}
  \caption{Seis cuadros obtenidos mediante \texttt{env.render()} al ejecutar la media
  determinista de REINFORCE en la semilla 1042. El episodio duró 160 pasos y obtuvo
  $-101.79$, por lo que la visualización también muestra que el entrenamiento no produjo
  una política consistentemente resuelta.}
  \label{fig:render}
\end{figure}

\clearpage

% =============================================================
\section{Task 3}
% =============================================================

\subsection{Contraste de las predicciones del Task 1}

\begin{respuesta}
\begin{center}
\begin{tabular}{p{0.09\textwidth}p{0.13\textwidth}p{0.68\textwidth}}
\toprule
\textbf{Pred.} & \textbf{Veredicto} & \textbf{Evidencia y explicación} \\
\midrule
P1 & \textbf{Refutada} & Actor-Critic no cruzó una media móvil de cero; REINFORCE
sí lo hizo en el episodio 778 y terminó $162.36$ puntos por encima al comparar los
promedios de los últimos 100. El mayor número de actualizaciones de TD(0) no compensó
el sesgo de un Critic impreciso ni la interferencia entre las pérdidas en el cuerpo
compartido. La caída de Actor-Critic hasta medias cercanas a $-1050$ alrededor del
episodio 500 muestra que más actualizaciones también amplifican un objetivo sesgado. \\
\addlinespace
P2 & \textbf{Confirmada} & La mediana de la norma fue $2287.58$ para REINFORCE y
$38.76$ para Actor-Critic; sus máximos fueron $30709.21$ y $515.61$. La comparación es
cualitativa porque una norma agrega todo un episodio y la otra promedia pasos, pero tanto
la diferencia de escala como la dispersión confirman la dirección predicha. \\
\addlinespace
P3 & \textbf{Refutada} & REINFORCE pasó de $1.4516$ a $1.4369$ nats de forma no
monótona. Actor-Critic descendió inicialmente hasta $1.3388$ en el episodio 206, pero
después aumentó hasta $1.8816$; no se volvió más determinista ni cayó más rápido de forma
sostenida. Los errores TD predominantemente negativos favorecieron reducir la
probabilidad de las muestras tomadas, lo cual puede lograrse aumentando $\sigma$. \\
\addlinespace
P4 & \textbf{Inconclusa} & La sección 1.4 anticipó que una penalización local de
suavidad favorecería Actor-Critic. Esta corrida no agregó
$-\lambda\lVert a_t-a_{t-1}\rVert^2$ ni modeló la dinámica del exoesqueleto, por lo que
esa afirmación no puede probarse con LunarLander. \\
\bottomrule
\end{tabular}
\end{center}
\end{respuesta}

\subsection{Norma del gradiente y pasos grandes}

\begin{respuesta}
Se definió una norma anormalmente alta mediante el criterio de Tukey
$Q_3+1.5\operatorname{IQR}$. Los umbrales resultaron $8739.39$ para REINFORCE y
$229.76$ para Actor-Critic, con 72 y 19 episodios atípicos, respectivamente. En
REINFORCE, el máximo ocurrió en el episodio 33: norma $30709.21$ y recompensa
$-588.87$. Entre los 64 episodios atípicos posteriores al episodio 20, 56 (87.5\%)
quedaron por debajo de la recompensa media de los 20 episodios previos. La recompensa
media en episodios atípicos fue $-321.88$, frente a $-77.89$ en los demás, y la
correlación de Pearson entre $\log(1+\lVert\nabla L_A\rVert)$ y recompensa fue
$-0.404$. Hay, por tanto, evidencia de que gradientes grandes coinciden con desempeño
peor, aunque una correlación observacional no demuestra causalidad.

Actor-Critic tuvo gradientes mucho menores, pero no estuvo libre de saltos: el máximo
del episodio 518 coincidió con recompensa $-1139.64$. En este método solo 8 de 18
atípicos evaluables cayeron por debajo de su media previa; además, el promedio por
episodio puede ocultar picos de una sola transición. El caso observado ilustra el
problema que motivó PPO: una actualización sin restricción puede cambiar demasiado la
razón $\pi_{\theta_{new}}(a\mid s)/\pi_{\theta_{old}}(a\mid s)$ y destruir una política
útil. El objetivo recortado de PPO no elimina gradientes ruidosos, pero impide que las
muestras sigan incentivando cambios una vez que la razón sale de $[1-\epsilon,
1+\epsilon]$, reduciendo el efecto de esos pasos.
\end{respuesta}

\subsection{Entropía de la política}

\begin{respuesta}
La hipótesis de una disminución gradual no se cumplió. REINFORCE permaneció cerca de su
valor inicial teórico: alcanzó un máximo de $1.4820$ en el episodio 174, un mínimo de
$1.4207$ en el 685 y terminó en $1.4369$. Su caída más fuerte en una ventana de 20
episodios fue de $0.0233$ nats entre los episodios 897 y 917, pero antes y después hubo
recuperaciones; no existe colapso sostenido de exploración.

Actor-Critic sí cayó rápidamente al inicio: su mayor descenso en 20 episodios fue
$0.0786$ entre los episodios 13 y 33. Tras el mínimo de $1.3388$ en el episodio 206,
la tendencia se invirtió y la entropía terminó en su máximo, $1.8816$. Agregar a la
pérdida un bono $-\beta S[\pi_\theta]$, como en PPO, mantendría exploración y protegería
contra un colapso prematuro de $\sigma$; sin embargo, en este Actor-Critic aumentaría una
dispersión que ya es excesiva. Sería necesario ajustar $\beta$, usar decaimiento o una
entropía objetivo, no asumir que más entropía siempre mejora el control.
\end{respuesta}

\subsection{Revisión bibliográfica}

\begin{respuesta}
Se revisó el trabajo de Luo et al., \enquote{Robust walking control of a lower limb
rehabilitation exoskeleton coupled with a musculoskeletal model via deep reinforcement
learning}, publicado en \emph{Journal of NeuroEngineering and Rehabilitation} en 2023
\cite{luo2023}. El problema es controlar un exoesqueleto de rehabilitación de miembro
inferior con ocho grados de libertad actuados ante interacciones humano--robot inciertas:
parálisis, debilidad muscular, condición sana y hemiparesia. Los autores entrenan en
simulación un controlador PPO con política estocástica, una red para la interacción
humana y una red supervisada de coordinación de 284 unidades musculotendinosas. PPO se
elige porque el objetivo recortado con $\epsilon=0.2$ limita cambios destructivos de la
política en un sistema continuo con contactos. Además, aleatorizan masa, inercia,
fricción, fuerza muscular y retardos para aumentar robustez.

En pruebas virtuales, el error angular RMSE quedó entre $2.048^\circ$ y $5.294^\circ$ en
las cuatro condiciones; el centro de presión permaneció dentro de la región estable de
$7\times11$ cm y la mayoría de torques quedó bajo 100 Nm. Una política entrenada solo
con humano pasivo obtuvo menos de 40\% de éxito con debilidad y 0\% con condición sana o
hemiparesia, lo que respalda la aleatorización de fuerza. El trabajo aborda directamente
problemas observados aquí: usa el recorte de PPO para los pasos grandes, una ventaja con
red de valor para reducir varianza, una recompensa de segunda diferencia de acciones,
un filtro pasa-bajos y límites de torque para suavidad. Persisten el sesgo del Critic, la
dependencia de la recompensa y el alto costo muestral (aproximadamente 20 millones de
muestras). Sobre todo, la validación fue solo simulada: no demuestra seguridad,
comodidad ni beneficio clínico en pacientes y PPO tampoco ofrece una garantía formal de
estabilidad.

\medskip
\lead{Referencia verificable.} DOI:
\href{https://doi.org/10.1186/s12984-023-01147-2}{10.1186/s12984-023-01147-2}; texto
completo: \url{https://pmc.ncbi.nlm.nih.gov/articles/PMC10024861/}.
\end{respuesta}

\subsection{Dictamen para el equipo directivo}

\begin{respuesta}
\lead{Punto de partida.} Recomendamos la arquitectura Actor-Critic como punto de partida
de investigación, pero no el Actor-Critic simple entrenado en este laboratorio. La
evidencia es mixta y debe comunicarse sin ocultarla: REINFORCE logró el mejor desempeño
($14.73$ frente a $-147.63$ en los últimos 100 episodios), mientras Actor-Critic sufrió
sesgo e interferencia del tronco compartido. Sin embargo, REINFORCE presentó una mediana
de gradiente 59 veces mayor y requiere trayectorias completas; esas propiedades son
inadecuadas para asignar una señal local de suavidad y aprender con datos físicos caros.
Actor-Critic permite crédito temporal y es la base del PPO usado por Luo et al.; por
ello es la dirección escalable, siempre que se reemplace TD(0) simple por estimación de
ventaja estable y se separen o balanceen cuidadosamente Actor y Critic.

\medskip
\lead{Escalamiento a PPO.} Sí recomendamos migrar a PPO \emph{antes} de considerar
hardware, pero no pasar directamente de PPO a un paciente. El recorte responde a los
picos asociados con caídas de recompensa y el paper aporta evidencia simulada pertinente.
PPO no garantiza acciones suaves: deben mantenerse la penalización de diferencias y
\textit{jerk}, el historial de acciones en el estado, límites duros de tasa, filtrado y
saturación de torque. Después se requiere aleatorización de dominio, validación
\textit{hardware-in-the-loop}, un supervisor de seguridad independiente y pruebas
progresivas sin carga humana antes de cualquier protocolo clínico.
\end{respuesta}

% =============================================================
\section{Repositorio}

\begin{respuesta}
Código fuente y notebook:
\url{https://github.com/MarioBetancourt04/Aprendizaje-Por-Refuerzo---Laboratorio-6}
\end{respuesta}

% =============================================================
\section{Uso de IA generativa}

\begin{respuesta}
\lead{Task 1.} Se utilizó Claude (Anthropic) en los cuatro incisos. No hubo un prompt
único que produjera la respuesta; el diálogo se organizó en tres etapas: explicación del
marco teórico, contraste del grupo con las diapositivas y redacción conjunta en
\LaTeX{}. La herramienta apoyó la formalización del operador $\arg\max$, el gradiente
de la Gaussiana, el compromiso sesgo--varianza y la ampliación del estado necesaria para
una recompensa que depende de la acción anterior.

\medskip
\lead{Prompt inicial de Task 1.}
\begin{quote}\itshape
``Vamos con otro lab de aprendizaje por refuerzo, hagamos el formato.''
\end{quote}
Luego se pidió explicar primero cada concepto para que el grupo pudiera interpretarlo
antes de construir la resolución. Funcionó porque separó comprensión y redacción: las
predicciones falsables y el tratamiento de $\sigma$ aprendible se revisaron contra el
enunciado en vez de aceptar directamente una respuesta generada.

\medskip
\lead{Task 2.} Se utilizó OpenAI GPT-5.6 Sol como apoyo para comprender, revisar y
ejecutar la implementación. El prompt de trabajo se formuló así:
\begin{quote}\itshape
``Ayúdame a revisar y explicar paso a paso el código que deseo implementar para comparar
REINFORCE con línea base y Actor-Critic. Parte del código que ya escribimos en el
notebook y explícame cómo se relaciona cada actualización con las fórmulas del curso.
Indícame qué cambios hacen falta, cómo instalar \texttt{gymnasium[box2d]} y cómo ejecutar
el notebook localmente usando el entorno que se activa con \texttt{cv}. Después ayúdame
a comprobar que ambos métodos entrenen durante 1000 episodios y que las métricas y
gráficas solicitadas se generen correctamente.''
\end{quote}
Funcionó porque pidió explicar la relación entre teoría y código antes de modificarlo,
partió del avance existente y especificó tanto las dependencias como el entorno de
ejecución. La herramienta señaló el optimizador no importado, la ausencia del
entrenamiento de Actor-Critic y que la norma del tronco compartido mezclaba las dos
pérdidas. El grupo revisó cada corrección y ejecutó la corrida completa.

\medskip
\lead{Task 3.} Para interpretar los resultados se utilizó el siguiente prompt:
\begin{quote}\itshape
``Ayúdame a leer las gráficas y métricas obtenidas, explicando cómo contrastar cada
predicción de Task 1 y cómo identificar picos de gradiente, caídas de recompensa y cambios
de entropía. También oriéntame para buscar y comprender un paper reciente que aplique PPO
a exoesqueletos, verificando autores, resultados y DOI, para que podamos redactar nuestras
propias conclusiones y recomendaciones.''
\end{quote}
Funcionó porque solicitó interpretación y fuentes verificables, no una conclusión sin
evidencia. Toda afirmación numérica se volvió a calcular desde
\texttt{lab6\_metrics.npz}; los datos del paper se contrastaron con el DOI y el texto
completo en PubMed Central. El grupo conservó los resultados refutados y limitó las
conclusiones del paper a simulación.

\medskip
\lead{Verificación realizada.} Se revisaron la derivación de
$\nabla_\theta \ln \pi_\theta$ y la demostración de insesgadez de la línea base. El
notebook se ejecutó de principio a fin en el entorno local, las figuras se regeneraron
desde los historiales de 1000 episodios y el PDF final se compiló desde este archivo. La
referencia bibliográfica, su DOI y las cifras citadas se contrastaron con la fuente
primaria.
\end{respuesta}

\begin{thebibliography}{9}
\bibitem{luo2023}
S. Luo, G. Androwis, S. Adamovich, E. Nunez, H. Su y X. Zhou,
\enquote{Robust walking control of a lower limb rehabilitation exoskeleton coupled with
a musculoskeletal model via deep reinforcement learning},
\emph{Journal of NeuroEngineering and Rehabilitation}, vol. 20, art. 34, 2023.
\href{https://doi.org/10.1186/s12984-023-01147-2}{doi:10.1186/s12984-023-01147-2}.
\end{thebibliography}

\end{document}
